\chapter{GIỚI THIỆU}
\label{chuong1}


\section{Đặt vấn đề}

Ô nhiễm không khí là một trong những thách thức môi trường và sức khỏe cộng đồng cấp bách nhất mà Việt Nam đang đối mặt. Theo Báo cáo Chất lượng Không khí Thế giới năm 2024 của IQAir, Hà Nội ghi nhận nồng độ PM2.5 trung bình năm là 45,4 $\mu$g/m\textsuperscript{3} --- cao gấp hơn chín lần so với ngưỡng khuyến nghị 5 $\mu$g/m\textsuperscript{3} của Tổ chức Y tế Thế giới và gần gấp đôi tiêu chuẩn quốc gia của Việt Nam (QCVN 05:2013) là 25 $\mu$g/m\textsuperscript{3}. Bụi mịn (PM2.5), được định nghĩa là các hạt lơ lửng trong không khí có đường kính khí động học nhỏ hơn 2,5 micromet, là chất ô nhiễm không khí có mối liên hệ mạnh nhất với các tác động sức khỏe bất lợi, bao gồm bệnh tim mạch, bệnh hô hấp và tử vong sớm. Nghiên cứu Gánh nặng Bệnh tật Toàn cầu ước tính rằng phơi nhiễm PM2.5 ngoài trời góp phần gây ra hơn 4 triệu ca tử vong sớm trên toàn thế giới mỗi năm, trong đó Đông Nam Á nằm trong số các khu vực bị ảnh hưởng nặng nề nhất.

Bất chấp mức độ nghiêm trọng của vấn đề, hệ thống quan trắc PM2.5 mặt đất khả dụng của Việt Nam còn thưa thớt và phân bố không đồng đều về mặt địa lý. Nghiên cứu này sử dụng dữ liệu từ 40 trạm quan trắc tự động do Trung tâm Quan trắc Môi trường Việt Nam (CEM) vận hành, cung cấp đủ dữ liệu PM2.5 liên tục theo giờ trong giai đoạn nghiên cứu 2023--2025 (37 trạm sau kiểm soát chất lượng; xem Mục 3.3), tập trung chủ yếu ở vùng đồng bằng sông Hồng phía Bắc và giảm mạnh ở miền Trung và miền Nam. Mạng lưới quốc gia của Việt Nam có quy mô lớn hơn nhưng vẫn còn khiêm tốn --- quy hoạch tổng thể quan trắc môi trường của chính phủ đặt mục tiêu 201 trạm quan trắc không khí tự động vào năm 2030, tăng từ khoảng một trăm trạm hiện nay, và chỉ một phần trong số đó, chủ yếu ở các thành phố trực thuộc trung ương, đo PM2.5 thay vì chỉ đo các khí được quy định (SO\textsubscript{2}, NO\textsubscript{2}, CO, TSP) \cite{vietnamplus2024}. Để so sánh, các nghiên cứu báo cáo R\textsuperscript{2} trên 0,85 cho Trung Quốc, Hoa Kỳ và châu Âu được huấn luyện trên các mạng lưới quốc gia gồm hàng trăm đến hàng nghìn trạm --- riêng Trung Quốc đã vượt quá 1.500 trạm --- nhiều hơn từ một đến hai bậc so với số trạm huấn luyện khả dụng ở đây. Sự phủ sóng thưa thớt và không đồng đều này đồng nghĩa với việc đại đa số người dân Việt Nam --- đặc biệt là những người ở vùng nông thôn, các thành phố nhỏ và các khu công nghiệp nằm ngoài các trung tâm đô thị lớn --- thiếu khả năng tiếp cận thông tin chất lượng không khí tại địa phương, và như Chương 4 và Chương 5 chứng minh, chính khoảng cách từ một vị trí bất kỳ đến trạm quan trắc khả dụng gần nhất, chứ không phải tổng số trạm toàn quốc hay mật độ trạm trên đơn vị diện tích, mới là yếu tố ràng buộc cuối cùng đối với độ chính xác dự đoán không gian.

Viễn thám vệ tinh mang đến một giải pháp tiềm năng cho khoảng trống quan trắc này. Các thiết bị trên các vệ tinh quỹ đạo cực và vệ tinh địa tĩnh quan sát bầu khí quyển Trái Đất ở độ phân giải không gian 1--10 km, cung cấp các phép đo hàng ngày hoặc nhiều lần trong ngày về độ sâu quang học sol khí (AOD), nồng độ khí vết, và các đặc tính bề mặt đất có tương quan với PM2.5 ở mức mặt đất. Khi kết hợp với dữ liệu tái phân tích khí tượng và các thuật toán học máy, các quan trắc vệ tinh này, về nguyên tắc, có thể được sử dụng để ước tính nồng độ PM2.5 mặt đất tại các vị trí không có trạm quan trắc. Hướng tiếp cận này đã được áp dụng rộng rãi trong tài liệu khoa học, với nhiều nghiên cứu báo cáo độ chính xác dự đoán cao (R\textsuperscript{2} > 0,8) cho các khu vực như Trung Quốc, Hoa Kỳ và châu Âu. Câu hỏi trọng tâm mà luận văn này đặt ra là liệu hiệu suất tương tự có thể đạt được cho Việt Nam hay không, với mạng lưới quan trắc thưa thớt, khí hậu nhiệt đới và các nguồn phát thải đặc thù của nước ta.


\section{Các giải pháp hiện có và hạn chế}

Tài liệu khoa học về ước tính PM2.5 dựa trên vệ tinh đã phát triển nhanh chóng trong thập kỷ qua. Các nghiên cứu tại Trung Quốc, sử dụng mạng lưới hơn 1.500 trạm mặt đất, đã chứng minh giá trị R\textsuperscript{2} vượt quá 0,85 với các mô hình học máy như random forests, XGBoost, và mạng nơ-ron sâu, tích hợp các đặc trưng từ MODIS AOD, khí tượng ERA5, dữ liệu sử dụng đất, và đầu ra của mô hình vận chuyển hóa học (CTM). Kết quả tương tự cũng được báo cáo cho Hoa Kỳ và châu Âu, nơi các mạng lưới quan trắc dày đặc cung cấp dữ liệu huấn luyện dồi dào.

Tuy nhiên, một quan sát phương pháp luận quan trọng đã nổi lên trong tài liệu khoa học gần đây: ước tính PM2.5 từ vệ tinh thực chất bao gồm nhiều bài toán con khác nhau, và phương pháp đánh giá phải phù hợp với bài toán đang được giải quyết. Đại đa số các nghiên cứu đã công bố đánh giá mô hình bằng đánh giá chéo k-fold ngẫu nhiên --- một phương pháp phù hợp cho bài toán \emph{nội suy thời gian} (lấp khoảng trống dữ liệu tại các trạm đã biết), trong đó mô hình được phép học từ lịch sử quan trắc của mỗi trạm. Nhưng bài toán có ý nghĩa vận hành tại Việt Nam là \emph{ngoại suy không gian} --- ước tính PM2.5 tại các vị trí chưa từng được quan trắc --- và bài toán này đòi hỏi một phương pháp đánh giá khác: đánh giá bỏ-một-trạm-ra (LOSO), trong đó toàn bộ dữ liệu của trạm đánh giá bị giữ lại khỏi quá trình huấn luyện.

Các nghiên cứu gần đây đã định lượng khoảng cách giữa hai bài toán này. Kawano et al. \cite{kawano2025}, công bố trên tạp chí Science Advances, báo cáo R\textsuperscript{2} = 0,85 cho nội suy thời gian (đánh giá chéo ngẫu nhiên) nhưng chỉ 0,67 cho ngoại suy không gian (đánh giá chéo không gian) --- khoảng cách 0,18 phản ánh sự khác biệt cơ bản giữa hai bài toán. Meyer et al. \cite{meyer2018} ghi nhận khoảng cách lên đến 0,66 (R\textsuperscript{2} từ 0,90 xuống 0,24) trong bài toán dự đoán không gian-thời gian về nhiệt độ không khí --- được trích dẫn ở đây như minh chứng phương pháp luận tổng quát chứ không phải như tiền lệ riêng cho PM2.5. Ploton et al. \cite{ploton2020}, công bố trên Nature Communications, báo cáo R\textsuperscript{2} = 0,53 cho nội suy nhưng chỉ 0,14 cho ngoại suy không gian trong lập bản đồ sinh khối. Các khoảng cách này không phản ánh lỗi phương pháp mà phản ánh sự khác biệt bản chất giữa hai bài toán: biết ``đây là loại vị trí nào'' so với phải suy luận điều đó từ các biến quan sát từ xa. Phần lớn tài liệu khoa học PM2.5 báo cáo kết quả cho bài toán dễ hơn (nội suy thời gian), trong khi bài toán có ý nghĩa vận hành (ngoại suy không gian) ít được nghiên cứu hơn.

Đối với Việt Nam cụ thể, tài liệu khoa học đã công bố còn hạn chế. Ngo et al. \cite{ngo2023} phát triển một mô hình hiệu ứng hỗn hợp cho Việt Nam sử dụng 9 trạm quan trắc trong giai đoạn 2012--2020, báo cáo R\textsuperscript{2} = 0,75. Tuy nhiên, kết quả này được thu được bằng đánh giá chéo dựa trên mẫu, và số lượng trạm ít đặt ra câu hỏi về khả năng tổng quát hóa không gian. Một số nghiên cứu đã giải quyết bài toán dự đoán PM2.5 ở các nước lân cận --- Thái Lan, lưu vực sông Mekong --- nhưng chưa nghiên cứu nào áp dụng đánh giá chéo không gian nghiêm ngặt (bỏ-một-trạm-ra hoặc tương đương) cho bộ dữ liệu Việt Nam.

Ngoài ra, các sản phẩm PM2.5 toàn cầu được xây dựng từ dữ liệu vệ tinh, như GHAP (Global High-resolution Air Pollutants) và MERRA-2 (Modern-Era Retrospective analysis for Research and Applications, Version 2), ngày càng được sử dụng để đánh giá phơi nhiễm ở các khu vực thưa dữ liệu. Tuy nhiên, các sản phẩm này chưa được đánh giá một cách có hệ thống so với các phép đo mặt đất tại Việt Nam, và độ chính xác của chúng trong môi trường nhiệt đới, gió mùa với đặc trưng phát thải riêng biệt --- giao thông chủ yếu bằng xe máy, đốt sinh khối, công nghiệp quy mô nhỏ --- vẫn chưa được biết đến.


\section{Mục tiêu và hướng giải pháp}

Mục tiêu chính của luận văn này là xây dựng và đánh giá một cách nghiêm ngặt một mô hình học máy được tăng cường bằng vệ tinh để ước tính nồng độ PM2.5 mặt đất theo giờ trên toàn lãnh thổ Việt Nam, với trọng tâm là bài toán ngoại suy không gian --- dự đoán tại các vị trí chưa từng được quan trắc. Luận văn sử dụng đánh giá chéo bỏ-một-trạm-ra (LOSO) làm thước đo đánh giá chính, phương pháp phù hợp cho bài toán này, đồng thời báo cáo kết quả nội suy thời gian (đánh giá chéo ngẫu nhiên) --- kỳ vọng ban đầu từ tài liệu khoa học --- để đặt kết quả trong bối cảnh và định lượng khoảng cách giữa kỳ vọng và thực tế.

Các mục tiêu cụ thể được đề ra như sau. Thứ nhất, xây dựng một bộ đặc trưng toàn diện tích hợp nhiều nguồn dữ liệu vệ tinh và tái phân tích --- bao gồm MODIS MAIAC AOD, Himawari-8/9 AOD, khí vết TROPOMI, khí tượng ERA5, đầu ra mô hình vận chuyển hóa học GEOS-CF và MERRA-2, khí hậu PM2.5 từ GHAP, ánh sáng ban đêm, và mật độ xây dựng --- đồng thời đánh giá đóng góp của từng nguồn dữ liệu vào độ chính xác dự đoán. Thứ hai, định lượng khoảng cách giữa kỳ vọng ban đầu (nội suy thời gian, R\textsuperscript{2} khoảng 0,80) và bài toán trọng tâm (ngoại suy không gian, R\textsuperscript{2} khoảng 0,20) trên cùng một mô hình và bộ dữ liệu, xác định rằng thông tin về bản sắc trạm là nguồn kỹ năng chính. Thứ ba, nghiên cứu vai trò của phân tầng theo chế độ ô nhiễm trong hiệu suất mô hình và mô tả đặc điểm phụ thuộc vòng tròn phát sinh khi biến phân nhóm chính là biến mục tiêu dự đoán. Thứ tư, đánh giá mô hình trên một mạng lưới độc lập gồm các cảm biến chi phí thấp và một trạm tham chiếu chuẩn của Đại sứ quán Hoa Kỳ không được sử dụng trong quá trình huấn luyện. Thứ năm, đánh giá các sản phẩm PM2.5 toàn cầu hiện có so với các phép đo mặt đất tại Việt Nam và xác định mức độ phù hợp của chúng làm tiên nghiệm không gian.

Phương pháp mô hình hóa dựa trên XGBoost với chính quy hóa DART (Dropouts meet Multiple Additive Regression Trees), được lựa chọn nhờ hiệu suất mạnh trên dữ liệu bảng không gian địa lý và khả năng xử lý giá trị thiếu --- một yêu cầu thực tiễn khi AOD vệ tinh chỉ khả dụng cho khoảng 15\% số giờ do mây che phủ trong khí hậu nhiệt đới Việt Nam.


\section{Đóng góp của luận văn}

Luận văn này có bốn đóng góp chính:

(i) Luận văn trình bày đánh giá có hệ thống đầu tiên về dự đoán PM2.5 dựa trên vệ tinh cho Việt Nam từ kỳ vọng ban đầu đến bài toán trọng tâm, chứng minh rằng cùng một mô hình đạt R\textsuperscript{2} khoảng 0,80 cho nội suy thời gian nhưng chỉ khoảng 0,20 cho ngoại suy không gian --- khoảng cách phản ánh sự khác biệt cơ bản giữa kỳ vọng và thực tế vận hành, nhất quán với các phát hiện của Kawano et al. \cite{kawano2025}, Meyer et al. \cite{meyer2018}, và Ploton et al. \cite{ploton2020}, và được mở rộng sang bối cảnh nhiệt đới mạng lưới thưa thớt.

(ii) Luận văn cung cấp đánh giá có hệ thống về các sản phẩm mô hình vận chuyển hóa học cho Việt Nam, cho thấy GEOS-CF có chu kỳ ngày đêm PM2.5 đảo ngược so với quan trắc mặt đất kèm theo sai số trung bình vượt quá 200\%, MERRA-2 --- mặc dù xấp xỉ không thiên lệch về trung bình --- chỉ đạt Chỉ số Phù hợp khoảng 0,39, thấp nhất trong tất cả các khu vực toàn cầu, và GHAP chỉ đạt khoảng 50\% mức phù hợp với phân loại mức ô nhiễm theo trạm --- chứng minh rằng các sản phẩm toàn cầu được hiệu chuẩn chủ yếu trên dữ liệu Đông Á và phương Tây hoạt động kém trong môi trường nhiệt đới Đông Nam Á.

(iii) Luận văn xác định, mô tả đặc trưng, và giải quyết phụ thuộc vòng tròn cơ bản trong dự đoán phân tầng theo chế độ ô nhiễm: phân tầng các trạm theo mức PM2.5 cải thiện R\textsuperscript{2} của mô hình khoảng 0,21, nhưng đòi hỏi phải biết trước biến được dự đoán. Thông qua kiểm tra có hệ thống nhiều phương pháp xấp xỉ --- cổng mềm, đường cơ sở dự đoán, dấu vân tay chế độ từ vệ tinh, gán dựa trên GHAP, phân cụm phát thải --- luận văn chứng minh rằng không có đại diện quan sát được nào từ vệ tinh giải quyết đáng tin cậy phụ thuộc này, với tất cả các phương pháp triển khai được dựa trên quan sát đều hội tụ ở R\textsuperscript{2} trung vị theo trạm khoảng 0,04. Sau đó, luận văn chỉ ra rằng phụ thuộc này vẫn có thể giải quyết được từ mạng lưới quan trắc chứ không phải từ các biến quan sát: mô hình định tuyến tiên nghiệm không gian neo đường cơ sở của mỗi trạm vào trung bình quan sát có trọng số khoảng cách của các trạm lân cận, phục hồi R\textsuperscript{2} trung bình theo trạm triển khai được là 0,197, so với ngưỡng trần oracle 0,203, không có phân loại sai nguy hiểm thấp-thành-cao, và tổng quát hóa tốt đến các cảm biến chi phí thấp chưa từng thấy như một bản đồ sàng lọc an toàn. Việc giải quyết này phụ thuộc vào khoảng cách giữa các trạm, khiến kỹ năng triển khai được đạt được trở thành hàm của mật độ mạng lưới chứ không phải một giới hạn trần duy nhất.

(iv) Luận văn đánh giá mô hình trên 39 trạm cảm biến chi phí thấp độc lập và một trạm tham chiếu chuẩn Đại sứ quán Hoa Kỳ không được sử dụng trong quá trình phát triển mô hình, đạt R\textsuperscript{2} = 0,68 tại Đại sứ quán và R\textsuperscript{2} trung vị = 0,53 trên các trạm LCS trong vùng phủ mạng lưới dày đặc --- kết quả cạnh tranh với các kết quả đánh giá chéo không gian tốt nhất được công bố trên quốc tế và chứng minh rằng khả năng dự đoán không gian của mô hình chủ yếu bị giới hạn bởi mật độ trạm. Các con số bên ngoài này được tính từ các bản ghi cảm biến chi phí thấp chỉ trong mùa khô dưới thiết lập chuyển đặc trưng từ trạm gần nhất (các biến dự đoán vệ tinh và khí tượng được lấy từ trạm quản lý gần nhất thay vì đo tại vị trí cảm biến), và nên được hiểu là kết quả chuyển đặc trưng và phù hợp cảm biến chứ không phải dự đoán không gian thực sự tại vị trí cảm biến.


\section{Cấu trúc luận văn}

Phần còn lại của luận văn này được tổ chức như sau.

Chương 2 trình bày cơ sở lý thuyết cần thiết để hiểu bài toán ước tính PM2.5 dựa trên vệ tinh. Chương xuất phát từ kỳ vọng ban đầu của tài liệu khoa học (nội suy thời gian) và cho thấy bài toán vận hành thực sự --- ngoại suy không gian --- là khó hơn một cách cơ bản. Chương tổng quan cơ sở vật lý của mối quan hệ AOD-PM2.5 trong môi trường nhiệt đới, khảo sát khung mô hình hóa XGBoost/DART, và trình bày các phương pháp đánh giá tương ứng.

Chương 3 mô tả phương pháp đề xuất một cách chi tiết. Chương bao gồm khu vực nghiên cứu và mạng lưới quan trắc --- 40 trạm quản lý, 83 cảm biến chi phí thấp, và một trạm Đại sứ quán Hoa Kỳ --- cùng các bộ dữ liệu vệ tinh và tái phân tích đa nguồn được sử dụng. Chương trình bày quy trình kỹ thuật đặc trưng tạo ra các biến dự đoán, quy trình kiểm soát chất lượng dữ liệu đã xác định và loại bỏ các cảm biến bị trục trặc, sơ đồ phân tầng theo mức ô nhiễm, và các phương pháp triển khai được phát triển để giải quyết phụ thuộc vòng tròn. Chương cũng mô tả mô hình vùng áp dụng cho mạng lưới dày đặc ở đồng bằng sông Hồng.

Chương 4 cung cấp phân tích chi tiết về các kết quả thí nghiệm. Chương xem xét các mẫu tầm quan trọng đặc trưng theo các tầng ô nhiễm, cho thấy nội suy không gian từ các trạm lân cận chi phối sức dự đoán của mô hình trong khi dấu vân tay chế độ từ vệ tinh cung cấp thông tin bổ sung. Chương đánh giá hiệu suất của các sản phẩm mô hình vận chuyển hóa học cho Việt Nam và trình bày nghiên cứu loại bỏ chỉ sử dụng vệ tinh. Chương đánh giá các bản đồ PM2.5 toàn cầu hiện có so với dữ liệu thực địa Việt Nam và phân tích mật độ trạm như ràng buộc then chốt đối với độ chính xác dự đoán không gian.

Chương 5 trình bày đánh giá thực nghiệm, xuất phát từ kỳ vọng ban đầu (đánh giá chéo ngẫu nhiên) và đi sâu vào bài toán trọng tâm (LOSO), cùng đánh giá độc lập trên cảm biến chi phí thấp. Chương trình bày kết quả phân tầng theo mức ô nhiễm, so sánh tất cả các cấu hình mô hình triển khai được, và ghi nhận toàn bộ các thí nghiệm đã khảo sát.

Chương 6 kết thúc luận văn với tổng hợp các phát hiện, khuyến nghị thực tiễn cho chiến lược quan trắc chất lượng không khí của Việt Nam, và các hướng phát triển trong tương lai.
